% Part: many-valued-logic
% Chapter: syntax-and-semantics
% Section: introduction

\documentclass[../../../include/open-logic-section]{subfiles}

\begin{document}

\olfileid{mvl}{syn}{int}

\olsection{परिचय}

अभिजात तर्कशास्त्रात $\lnot$, $\land$, $\lor$, $\lif$ आणि
$\liff$ ही विधानीय संयोजके वापरून !!{propositional variable}s पासून
घडवलेल्या !!{formula}s चा विचार करतो. अभिजात तर्कशास्त्राची
चिन्हार्थमीमांसा परिभाषित करताना आपण $\True$ आणि $\False$ ही दोन
सत्यतामूल्ये वापरतो. !!a{valuation}~$\pAssign{v}$ मध्ये
!!{propositional variable}s ना अर्थ देतो; ते !!{propositional variable}s ना $\True$ आणि
$\False$ ही सत्यतामूल्ये नेमून देते. मग कोणतेही !!{valuation}
कोणत्याही !!{formula}~$!A$ चे $\pValue{v}(!A)$ हे सत्यतामूल्य
ठरवते; आणि !!^a{formula} ची !!a{valuation}~$\pAssign{v}$ मध्ये
पूर्ति होते, $\pSat{v}{!A}$, जर आणि फक्त जर
$\pValue{v}(!A) = \True$.

बहुमूल्य तर्कशास्त्रे $\True$ आणि $\False$ यांपलीकडील
सत्यतामूल्यांना परवानगी देऊन अभिजात द्विमूल्य तर्कशास्त्राचे
सामान्यीकरण करतात. म्हणून बहुमूल्य तर्कशास्त्रात
!!a{valuation}~$\pAssign{v}$ हे प्रत्येक !!{propositional variable}~$p$
ला संभाव्य सत्यतामूल्यांपैकी एक मूल्य नेमून देणारे फलन असते.
अनुमत सत्यतामूल्यांच्या संचाला आपण साधारणपणे~$V$ म्हणू.
$V = \{\True, \False\}$ असलेले अभिजात तर्कशास्त्र हेही
बहुमूल्य तर्कशास्त्र आहे; त्यात $\pValue{v}(!A)$ हे सत्यतामूल्य
संयोजकांच्या परिचित वैशिष्ट्यपूर्ण सत्यताकोष्टकांनुसार ठरते.

अधिक सत्यतामूल्ये समाविष्ट केल्यावर ``आणि,'' ``किंवा,'' ``नाही,''
आणि ``जर---तर'' असे वाचत असलेल्या संयोजकांसाठी
$\pValue{v}(!A)$ कसे काढायचे, याचे एकाहून अधिक स्वाभाविक
पर्याय असतात. म्हणून बहुमूल्य तर्कशास्त्र केवळ सत्यतामूल्यांच्या
संचानेच नव्हे, तर प्रत्येक संयोजकासाठी निवडलेल्या
\emph{सत्यता-फलनांनी}ही निश्चित होते. तथापि, तर्कशास्त्र~$\Log L$
साठी ही निवड झाल्यावर, अभिजात तर्कशास्त्राप्रमाणेच
$\pValue{v}(!A)[\Log L]$ हे सत्यतामूल्य त्या सत्य-मूल्यांकनाने
एकमेव रीतीने ठरते. अभिजात तर्कशास्त्राप्रमाणे बहुमूल्य
तर्कशास्त्रेही \emph{सत्यता-फलनात्मक} आहेत.

हे चिन्हार्थविषयक पायाभूत घटक मिळाल्यावर सर्वतः सत्यता,
तार्किक निष्पन्नता आणि पूर्ततायोग्यता यांच्या समकक्ष
चिन्हार्थविषयक संकल्पना परिभाषित करता येतात. अभिजात
तर्कशास्त्रात !!a{formula} सर्वतः सत्य असते, जर
कोणत्याही~$\pAssign{v}$ साठी त्याचे सत्यतामूल्य
$\pValue{v}(!A) = \True$ असेल. बहुमूल्य तर्कशास्त्रात ही
संकल्पना थोडी व्यापक करावी लागते. प्रथम, सत्यतामूल्यांच्या~$V$
संचात $\True$ असणे आवश्यक नाही. उदाहरणार्थ,
काही बहुमूल्य तर्कशास्त्रे $0$ आणि~$1$ यांदरम्यानच्या सर्व
परिमेय संख्या सत्यतामूल्ये म्हणून वापरतात. अशा वेळी $1$
सहसा $\True$ ची भूमिका बजावते. इतर तर्कशास्त्रांत
एकच नव्हे, तर अनेक सत्यतामूल्ये ती भूमिका बजावतात.
म्हणून प्रत्येक बहुमूल्य तर्कशास्त्रात \emph{निर्दिष्ट
सत्यतामूल्यां}चा संच~$V^+$ असावा अशी अट घालतो.
मग !!a{formula} ची !!a{valuation}~$\pAssign{v}$ मध्ये
पूर्ति होते, $\pSat{v}{!A}[\Log L]$, जर आणि फक्त जर
$\pValue{v}(!A)[\Log L] \in V^+$. !!^a{formula}~$!A$
हे त्या तर्कशास्त्रातील सर्वतः सत्य सूत्र असते,
$\Entails[\Log L] !A$, जर आणि फक्त जर
कोणत्याही~$\pAssign{v}$ साठी $\pValue{v}(!A) \in V^+$.
शेवटी, !!{formula}s च्या संचापासून $!A$ निष्पन्न होते,
$\Gamma \Entails[\Log L] !A$, जर~$\Gamma$ मधील सर्व
!!{formula}s ची पूर्ति करणारे प्रत्येक !!{valuation}~$!A$
चीही पूर्ति करत असेल.

\end{document}
